\documentclass[17pt,a4paper,landscape]{extarticle}
\usepackage[margin=2.35cm,top=0.12cm,left=1.05cm]{geometry}
\usepackage{amsmath}
\usepackage{amssymb}
\usepackage{tasks}
\usepackage{tikz}
\usepackage{xcolor}
\usepackage[sfdefault,lf]{FiraSans}
\renewcommand{\familydefault}{\sfdefault}
\setlength{\parindent}{0pt}
\pagestyle{empty}
\settasks{label=\textbf{\arabic*.}, label-width=2.2em, item-indent=2.95em, column-sep=2.1cm, after-item-skip=4.7em}
\renewcommand{\arraystretch}{1.15}
\everymath{\displaystyle}

\begin{document}
\sffamily
\boldmath

\boldmath

{\LARGE \textbf{Excellence}}
\begin{tasks}(3)
\task A student says \mbox{$0.35=35\%=\frac{35}{100}=\frac{7}{20}$}. Are they correct? Explain.
\task A student says \mbox{$\frac{3}{8}=0.38=38\%$}. Are they correct? Explain.
\task Convert \mbox{$87.5\%$} to a fraction in simplest form.
\task Convert \mbox{$\frac{13}{20}$} to a decimal and a percentage.
\task Convert \mbox{$0.625$} to a fraction in simplest form and a percentage.
\task Which is greater: \mbox{$65\%$} or \mbox{$\frac{2}{3}$}? Show enough working to justify.
\task Which is smaller: \mbox{$0.48$} or \mbox{$\frac{1}{2}$}? Explain.
\task Fill in the blank: \mbox{$\frac{\square}{8}=0.625=62.5\%$}.
\task Fill in the blank: \mbox{$45\%=0.45=\frac{9}{\square}$}.
\task Fill in the blank: \mbox{$0.2=20\%=\frac{\square}{\square}$} in simplest form.
\task Put in descending order: \mbox{$\frac{7}{10},\ 68\%,\ 0.72$}.
\task Put in ascending order: \mbox{$12.5\%,\ \frac{1}{5},\ 0.18$}.
\task Which does not belong: \mbox{$0.4,\ 40\%,\ \frac{4}{10},\ \frac{1}{5}$}?
\task Complete: a number equivalent to \mbox{$\frac{3}{5}$} is \rule{2cm}{0.15mm} as a percentage.
\task Complete: a decimal equivalent to \mbox{$125\%$} is \rule{2cm}{0.15mm}.
\task A test score was \mbox{$18$} out of \mbox{$24$}. Write this as a fraction in simplest form, a decimal, and a percentage.
\task A bottle is \mbox{$0.75$} full. Write this as a fraction and a percentage.
\task Explain why \mbox{$\frac{1}{3}$} cannot be written as a terminating decimal even though it can be written as a percentage.
\end{tasks}

\end{document}
